\documentclass[11pt]{article}

\usepackage[utf8]{inputenc}
\usepackage[T1]{fontenc}
\usepackage[brazil]{babel}
\usepackage[expansion=false]{microtype}

\usepackage[a4paper,margin=2.5cm]{geometry}
\usepackage{amsmath,amssymb}
\usepackage{booktabs}
\usepackage{graphicx}
\usepackage{float}
\usepackage{xcolor}
\usepackage{enumitem}
\usepackage[hidelinks]{hyperref}
\usepackage{caption}
\captionsetup{font=small,labelfont=bf}
\usepackage{abstract}
\usepackage{authblk}

\title{\textbf{Cosseno 0,99 não é compreensão:\\
\large anisotropia, polaridade e a competência jurídica\\
diante da recuperação densa}}

\author[ ]{Pedro Borges Mourão}
\affil[ ]{\small Promotor de Justiça, Mestre em Direito, Especialista em Segurança da Informação\\
e pesquisador na interseção entre Direito, inteligência artificial e forense digital.}

\date{\today}

\begin{document}
\maketitle

% ======================================================================
\begin{abstract}
\noindent
Circula entre operadores do Direito a suposição de que dominar uma ferramenta de
inteligência artificial é uma questão de aprender comandos. Este artigo argumenta
que essa suposição é falsa e que seu custo não é a ineficiência, mas a confiança
indevida em sistemas cujo modo de falha é invisível para quem não entende como
eles representam a linguagem. O componente que examinamos é a recuperação densa
por \textit{sentence embeddings}, base dos sistemas de busca jurídica e dos
\textit{pipelines} de geração aumentada por recuperação (RAG). Construímos um
conjunto de 40 pares mínimos forenses em português e medimos quatro modelos
abertos, de três modelos consolidados (MiniLM, mpnet, distiluse) a um modelo de
fronteira (multilingual-e5-large). O achado central é contraintuitivo e tem
consequência prática direta. Os modelos preservam a polaridade --- a distinção
entre afirmação e negação, entre polo ativo e passivo --- como um sinal
estruturalmente fraco, e a fraqueza é mais severa não na negação adverbial, mas
na inversão de polo processual (autor\,$\leftrightarrow$\,réu). Esse padrão
\emph{atravessa as gerações de modelos}: ele aparece nos três modelos antigos e
reaparece, idêntico, no modelo de fronteira quando se corrige a geometria do
espaço. O modelo de fronteira, porém, opera em um regime fortemente
\emph{anisotrópico} --- a similaridade de cosseno entre duas sentenças jurídicas
quaisquer, ainda que sem relação, é de $0{,}84$ --- de modo que ele atribui a uma
sentença e à sua negação um cosseno bruto de $0{,}99$. Normalizada pelo
\textit{baseline} anisotrópico, a discriminação reaparece e é, em termos
relativos, \emph{menor} que a dos modelos antigos ($\textit{gap}$ normalizado de
$0{,}09$ contra $0{,}27$); no estrato mais severo, a inversão de polo
(autor\,$\leftrightarrow$\,réu), o \textit{gap} chega a ser \emph{negativo} --- o
e5 considera a tese e sua inversão mais próximas entre si do que duas paráfrases
legítimas. Na tarefa de recuperação, o modelo de fronteira acerta
\emph{mais} ($10/10$ contra $8/10$), porém por uma margem de cosseno cujo mínimo
é de cinco milésimos. A conclusão inverte a intuição corrente: o progresso da
engenharia de \textit{embeddings} não tornou a polaridade mais robusta, mas
deslocou sua fragilidade para um regime onde o operador, lendo um cosseno de
$0{,}99$, conclui certeza máxima exatamente onde a discriminação é mais precária.
A modernização do modelo disfarçou a armadilha. Reconhecê-la exige uma
alfabetização linguística e geométrica que a formação jurídica não fornece. Todo
o material é reproduzível.

\vspace{0.6em}
\noindent\textbf{Palavras-chave:} inteligência artificial jurídica; recuperação
densa; \textit{embeddings}; anisotropia; polaridade; geração aumentada por
recuperação; prova digital.
\end{abstract}

% ======================================================================
\section{Introdução: uma suposição confortável e cara}

Há uma suposição confortável circulando nos escritórios e órgãos públicos que
operam a justiça no Brasil: a de que dominar uma ferramenta de inteligência
artificial é uma questão de aprender comandos, como quem aprende os atalhos de
um editor de texto. Descobre-se o botão certo, memorizam-se algumas fórmulas de
prompt e, pronto, está-se operando. Essa suposição é falsa, e o custo de
mantê-la não é a ineficiência. É a confiança indevida em um sistema cujo modo de
falha o operador não enxerga.

O argumento deste texto é específico, e convém delimitá-lo com precisão para que
não seja confundido com uma crítica genérica à IA. Não se defende que advogados
aprendam técnica de prompt. O objeto que examinamos é um componente preciso da
infraestrutura de IA jurídica: a \emph{recuperação densa por
\textit{embeddings}}, isto é, o mecanismo que converte peças e consultas em
vetores numéricos e os compara por similaridade. Esse componente é o coração dos
sistemas de busca jurisprudencial modernos e a primeira etapa dos
\textit{pipelines} de geração aumentada por recuperação (RAG), que alimentam o
modelo gerativo com os documentos que ele usará para responder. Se essa etapa
entrega ao gerador o documento de sentido invertido --- o precedente que
contradiz a tese, em vez do que a sustenta ---, todo o restante do sistema opera
sobre material corrompido na origem. É essa etapa, e não a geração de texto em si,
que este artigo diagnostica.

A contribuição é dupla. Primeiro, ela mede empiricamente, em português jurídico,
como esse componente trata a polaridade, e mostra que a fragilidade observada não
é defeito de um modelo específico, mas propriedade estrutural que atravessa
gerações. Segundo, e é o achado que reorganiza o artigo, ela demonstra que o
progresso tecnológico não eliminou essa fragilidade: deslocou-a para um regime
geométrico onde ela se torna mais difícil de perceber. O operador que conclui
``máxima confiança'' diante de um escore de similaridade altíssimo está, nesse
regime, cometendo precisamente o erro que a competência jurídica não o ensina a
evitar.

% ======================================================================
\section{O caso que expõe a natureza do problema}
\label{sec:caso}

Tomem-se duas frases que dizem o oposto uma da outra: ``o réu compareceu à
audiência'' e ``o réu não compareceu à audiência''. Para um operador do Direito,
essas frases são antípodas. A presença ou ausência do réu muda o procedimento,
dispara ou não a revelia (art.~344 do CPC), pode redesenhar o ônus da prova. A
diferença entre elas é a diferença entre dois processos.

Submetam-se essas duas frases a um modelo de \textit{embedding}. A similaridade
que o modelo atribui ao par é alta. A causa do fenômeno precede qualquer
arquitetura específica: a representação distribucional do significado, formulada
por Zellig Harris em 1954, segundo a qual o significado de uma palavra é capturado
pelos contextos em que ela aparece \cite{harris1954}. Modelos de \textit{embedding}
aprendem que ``compareceu'' e ``não compareceu'' significam coisas parecidas
porque aparecem em contextos textuais quase idênticos. O advérbio de negação é,
para a estatística distribucional, uma perturbação mínima no contexto. Para o
Direito, é a inversão completa do sentido.

A questão empírica que organiza o restante do artigo tem três camadas, e
mantê-las separadas é a principal exigência metodológica. A primeira: \emph{quão}
alta é essa similaridade, comparada à de uma paráfrase legítima, quando se isola a
sobreposição lexical? A segunda: esse comportamento muda com a geração do modelo,
ou modelos de fronteira o corrigem? A terceira, que se revelou decisiva: a
similaridade absoluta sequer é a métrica certa, dado que o espaço vetorial pode
ser geometricamente enviesado? As três camadas correspondem às três contribuições
empíricas das Seções~\ref{sec:resultados} e \ref{sec:aniso}.

% ======================================================================
\section{Método}
\label{sec:metodo}

\subsection{Pares mínimos e o problema do controle}

Um \emph{par mínimo} compartilha quase toda a superfície lexical e difere em um
ponto que inverte o sentido. Construímos 40 pares em português jurídico,
estratificados pelo \emph{mecanismo} de inversão (Tabela~\ref{tab:estratos}).

\begin{table}[H]
\centering
\small
\begin{tabular}{lll}
\toprule
Estrato & $n$ & Exemplo (frase A / frase B) \\
\midrule
Negação adverbial & 6 & \emph{O réu compareceu\ldots} / \emph{O réu não compareceu\ldots} \\
Antonímia de núcleo & 6 & \emph{\ldots foi deferido} / \emph{\ldots foi indeferido} \\
Inversão de polo & 16 & \emph{O autor deve indenizar o réu} / \emph{O réu deve indenizar o autor} \\
Quantificador/modal & 4 & \emph{Todos os requisitos\ldots} / \emph{Nenhum dos requisitos\ldots} \\
Controle pareado & 8 & \emph{\ldots foi indeferido} / \emph{\ldots foi rejeitado} \\
\bottomrule
\end{tabular}
\caption{Estratos do conjunto. O estrato de \emph{inversão de polo} foi ampliado
para $n=16$ por concentrar o efeito mais forte. Os oito pares de \emph{controle}
são paráfrases legítimas, lexicalmente pareadas aos contrastivos (Jaccard
$0{,}79$ vs $0{,}75$), de modo que a comparação difere apenas na variável de
interesse: preserva \emph{versus} inverte o sentido.}
\label{tab:estratos}
\end{table}

\subsection{Modelos, medida e as três camadas de análise}

Medimos quatro modelos abertos de \textit{sentence embedding} multilíngues:
três consolidados e amplamente usados em produção --- \texttt{paraphrase\-multilingual\-MiniLM\-L12\-v2}, \texttt{paraphrase\-multilingual\-mpnet\-base\-v2} e
\texttt{distiluse\-base\-multilingual\-cased\-v2} --- e um modelo de fronteira,
\texttt{multilingual\-e5\-large}, treinado com aprendizado contrastivo em larga
escala, da biblioteca \texttt{sentence-transformers} \cite{reimers2019}. O e5
exige prefixos \texttt{query:}/\texttt{passage:}; aplicamo-los de forma simétrica
(\texttt{query:} em ambas as sentenças) na medição intrínseca e assimétrica na
recuperação, conforme o treino do modelo.\footnote{A escolha do prefixo
\texttt{query:} para a medição intrínseca segue a instrução oficial dos autores do
e5, que recomendam esse prefixo para tarefas \emph{simétricas} como similaridade
semântica e recuperação de paráfrases, reservando o par \texttt{query:}/\texttt{passage:}
para tarefas assimétricas de recuperação de passagens longas. A anisotropia
observada não é, ademais, artefato da escolha de prefixo: os próprios autores
atribuem a concentração dos escores de cosseno na faixa de $0{,}7$ a $1{,}0$ ao
uso de temperatura baixa ($0{,}01$) na perda contrastiva InfoNCE, propriedade do
treino que independe do prefixo. O \textit{script} de replicação inclui a
verificação direta do \textit{baseline} sob ambos os protocolos.}

A análise tem três camadas. A camada \emph{intrínseca} mede o cosseno entre uma
sentença e sua negação. A camada \emph{geométrica} mede o \textit{baseline}
anisotrópico --- o cosseno médio entre sentenças jurídicas sem relação --- e
normaliza por ele. A camada \emph{extrínseca} mede, numa tarefa de recuperação,
se a sentença correta é rankeada acima da oposta. As três são necessárias porque,
como a Seção~\ref{sec:aniso} demonstra, a camada intrínseca isolada conduz a
conclusões erradas em espaços anisotrópicos. Toda comparação agregada acompanha
intervalo de confiança de 95\% por \textit{bootstrap} ($5{.}000$ reamostragens). A
execução é determinística e reproduzível.

% ======================================================================
\section{Resultados: a fragilidade da polaridade atravessa as gerações}
\label{sec:resultados}

A Tabela~\ref{tab:matriz} e a Figura~\ref{fig:heat} reportam o cosseno bruto por
modelo e estrato. Dois fatos se mantêm em todos os modelos. Primeiro, a negação
adverbial é o estrato \emph{menos} problemático --- a presença do marcador ``não''
dá ao modelo um sinal lexical a que se agarrar. Segundo, a inversão de polo é o
estrato \emph{mais} severo: trocar ``autor'' por ``réu'', dois \textit{tokens} de
altíssima frequência em contextos idênticos, produz o cosseno mais alto em todos
os quatro modelos, chegando a $0{,}993$ no e5.

\begin{table}[H]
\centering
\small
\begin{tabular}{lrrrrr}
\toprule
Estrato & MiniLM & mpnet & distiluse & e5 (bruto) & e5 (norm.) \\
\midrule
negacao\_adverbial & 0.550 & 0.490 & 0.844 & 0.958 & 0.743 \\
antonimia\_nucleo & 0.724 & 0.769 & 0.925 & 0.977 & 0.857 \\
inversao\_polo & 0.958 & 0.967 & 0.972 & 0.993 & 0.958 \\
quant\_modal & 0.693 & 0.562 & 0.901 & 0.975 & 0.846 \\
controle\_parafrase & 0.919 & 0.940 & 0.924 & 0.991 & 0.943 \\
\midrule
\textit{baseline aniso.} & 0.07 & $\sim$0.10 & $\sim$0.14 & \multicolumn{2}{c}{0.84} \\
\bottomrule
\end{tabular}

\caption{Cosseno bruto por estrato e modelo, com o \textit{baseline} anisotrópico
de cada espaço na última linha. Note o salto do \textit{baseline}: $0{,}07$ nos
modelos antigos, $0{,}84$ no e5.}
\label{tab:matriz}
\end{table}

\begin{figure}[H]
\centering
\includegraphics[width=0.92\textwidth]{v3_figA.pdf}
\caption{Cosseno bruto dos quatro modelos. O e5 esmaga todos os estratos perto de
$0{,}99$, inclusive os contrastivos. Lido ingenuamente, isso sugeriria que o e5 é
o \emph{pior} modelo --- conclusão que a Seção~\ref{sec:aniso} mostra ser um
artefato geométrico.}
\label{fig:heat}
\end{figure}

A leitura ingênua da Figura~\ref{fig:heat} seria que o e5, por dar cosseno mais
alto aos pares opostos, é o mais cego à polaridade. Essa leitura é falsa, e o erro
que ela comete é exatamente o erro que o operador jurídico comete diante da
ferramenta. Corrigi-lo exige tratar a geometria do espaço.

% ======================================================================
\section{A anisotropia muda o que o cosseno significa}
\label{sec:aniso}

\subsection{O fenômeno}

Em espaços vetoriais de alta dimensão produzidos por transformadores, os vetores
tendem a se concentrar em um cone estreito, em vez de se distribuírem pela esfera.
Essa propriedade, a \emph{anisotropia}, foi caracterizada por Ethayarajh
\cite{ethayarajh2019} e faz com que sentenças sem qualquer relação semântica
exibam cosseno alto. A consequência é direta: em um espaço anisotrópico, o valor
absoluto do cosseno deixa de ser interpretável, e apenas a posição \emph{relativa}
ao \textit{baseline} do espaço carrega informação. O fenômeno é especialmente
pronunciado em modelos multilíngues \cite{hammerl2023}.

Medimos o \textit{baseline} de cada modelo como o cosseno médio entre sentenças
jurídicas não relacionadas. O resultado separa as duas gerações de forma nítida:
o \textit{baseline} dos três modelos antigos fica entre $0{,}07$ e $0{,}14$ ---
espaços praticamente isotrópicos ---, ao passo que o do e5 é de $0{,}84$. No
espaço do e5, duas sentenças jurídicas quaisquer já partem de uma similaridade de
$0{,}84$; o cosseno de $0{,}99$ entre uma sentença e sua negação representa, sobre
essa base, um acréscimo pequeno. Dizer que o e5 ``vê as duas frases como $99\%$
idênticas'' é, portanto, um erro de leitura: ele vê quase tudo como muito similar,
e a pergunta correta é quanto a polaridade desloca o vetor \emph{acima} desse piso.

\subsection{Normalização e o reaparecimento da estrutura}

Para tornar os modelos comparáveis, normalizamos o cosseno pelo \textit{baseline}
de cada espaço, seguindo o procedimento de pares mínimos de Lin e Chen
\cite{heros2023}: a similaridade normalizada remapeia o \textit{baseline} para
zero, isolando o sinal semântico acima do ruído geométrico. A
Figura~\ref{fig:central} mostra o efeito no e5. À esquerda, o cosseno bruto:
indistinguível a olho nu, tudo entre $0{,}96$ e $0{,}99$. À direita, o cosseno
normalizado: a estrutura reaparece, e é a \emph{mesma} estrutura dos modelos
antigos --- a negação adverbial discriminada ($0{,}74$), a inversão de polo colada
no controle ($0{,}96$ contra $0{,}94$).

\begin{figure}[H]
\centering
\includegraphics[width=0.98\textwidth]{v3_figB.pdf}
\caption{O achado central. À esquerda, o que o operador vê: cosseno bruto do e5,
uniforme perto de $0{,}99$. À direita, o que de fato ocorre: removida a
anisotropia, a inversão de polo (azul) praticamente não se separa do controle de
paráfrase (verde). A fragilidade estava lá, oculta pela geometria.}
\label{fig:central}
\end{figure}

O dado que sintetiza o argumento é o \textit{gap} normalizado entre contrastivos e
controle. Nos modelos antigos, medido sobre o mesmo corpus jurídico, o \textit{gap}
do MiniLM é de $0{,}27$. No e5, é de $0{,}09$. Em termos relativos, descontada a
geometria, o modelo de fronteira discrimina a polaridade \emph{menos} que o modelo
antigo, não mais. O aprendizado contrastivo em larga escala, que a intuição
supõe ter ``resolvido'' a insensibilidade ao sentido, não a resolveu para a
inversão de polo; apenas a transferiu para um regime onde ela é mais difícil de
detectar. Esse $0{,}09$ é a média de todos os estratos contrastivos; quando se
isola a inversão de polo, o estrato mais severo, o \textit{gap} torna-se
\emph{negativo} (Seção~\ref{sec:limita}) --- o e5 chega a considerar a tese e sua
inversão mais próximas entre si do que duas paráfrases legítimas.

% ======================================================================
\section{O paradoxo: o modelo que acerta mais esconde melhor a fragilidade}
\label{sec:paradoxo}

Poder-se-ia objetar que tudo isso é geometria sem consequência, desde que a
recuperação funcione. Medimos, então, a tarefa-fim: para 10 consultas temáticas,
com a sentença correta, sua oposta e doze distratores no corpus, o modelo rankeia
a correta acima da oposta? A Figura~\ref{fig:paradoxo} reporta o resultado, e ele
é o ponto mais importante do artigo.

\begin{figure}[H]
\centering
\includegraphics[width=0.86\textwidth]{v3_figC.pdf}
\caption{O paradoxo do e5. Em azul, o acerto de \textit{ranking}: o e5 acerta
\emph{mais} ($10/10$) que os modelos antigos ($8/10$). Em vermelho, a margem
mínima de cosseno que sustenta esses acertos: no e5 ela é de cinco milésimos. O
modelo melhor na tarefa é o que opera mais perto da indistinção.}
\label{fig:paradoxo}
\end{figure}

O e5 é o melhor dos quatro na tarefa de recuperação: acerta todos os dez casos,
contra oito dos modelos antigos, e nunca rankeia a oposta acima da correta. Mas a
margem de cosseno que separa a correta da oposta é, no e5, de $0{,}005$ no pior
caso --- sete vezes menor que a do MiniLM. O modelo que mais acerta é o que decide
por margens mais finas, embutidas num espaço onde o cosseno absoluto de tudo é
$0{,}99$.

Esta é a armadilha, e é estritamente o erro que a competência jurídica não
previne. O operador que projeta ou contrata um sistema de busca, ao comparar
modelos, vê o e5 entregar similaridades de $0{,}99$ e conclui que ele ``entende''
os documentos com altíssima confiança. O que ele não vê é que essa confiança
aparente repousa sobre uma discriminação de polaridade de cinco milésimos, que é a
primeira coisa a colapsar sob qualquer perturbação realista --- um documento mais
longo que dilui o peso dos \textit{tokens} de polo, uma consulta ambígua, um
\textit{reranker} mal calibrado, ruído de OCR. O estudo de Stanford sobre
ferramentas jurídicas, no nível do produto, observa que documentos apenas
textualmente parecidos são recuperados como relevantes e degradam a resposta
\cite{magesh2025}; nossos resultados localizam, no componente de \textit{embedding}
e na sua geometria, um dos mecanismos desse fenômeno, e mostram que ele se agrava,
em vez de desaparecer, à medida que os modelos avançam --- porque o avanço vem
acompanhado de maior anisotropia.

% ======================================================================
\section{A inversão de polo resiste à atenção cruzada}
\label{sec:cross}

Poder-se-ia esperar que o problema da polaridade fosse específico do bi-encoder,
que comprime cada sentença em um vetor isolado antes de compará-las, e que uma
arquitetura que processasse o par \emph{conjuntamente} o resolvesse. Testamos essa
hipótese submetendo os 40 pares a um \textit{cross-encoder} de inferência de
linguagem natural, \texttt{nli-deberta-v3-base}, que recebe as duas sentenças ao
mesmo tempo, com atenção cruzada entre todos os \textit{tokens}, e as classifica
como \emph{contradição}, \emph{implicação} ou \emph{neutro}. Se a distinção de
polaridade fosse recuperável por atenção conjunta, esperaríamos que todos os pares
contrastivos fossem classificados como contradição e o controle de paráfrase, não.

A Tabela~\ref{tab:cross} e a Figura~\ref{fig:cross} mostram que essa expectativa se
confirma para a maioria dos estratos, mas falha exatamente onde o bi-encoder
também falha pior. A negação adverbial e os quantificadores são detectados como
contradição em $100\%$ dos casos ($6/6$ e $4/4$); a antonímia de núcleo, em
$5/6$. O controle de paráfrase é corretamente \emph{não} classificado como
contradição ($2/8$), validando o instrumento. A inversão de polo, porém, é
detectada como contradição em apenas $9$ de $16$ casos.

\begin{table}[H]
\centering
\small
\begin{tabular}{lrrrr}
\toprule
Estrato & Contradição & Implicação & Neutro & $n$ \\
\midrule
negacao\_adverbial & 6 & 0 & 0 & 6 \\
antonimia\_nucleo & 5 & 1 & 0 & 6 \\
inversao\_polo & 9 & 6 & 1 & 16 \\
quant\_modal & 4 & 0 & 0 & 4 \\
controle\_parafrase & 2 & 6 & 0 & 8 \\
\bottomrule
\end{tabular}

\caption{Classificação do \textit{cross-encoder} NLI por estrato. A inversão de
polo é a única em que a detecção de contradição falha sistematicamente, e o modo
de falha é revelador: 6 dos 16 pares são classificados como \emph{implicação}.}
\label{tab:cross}
\end{table}

\begin{figure}[H]
\centering
\includegraphics[width=0.9\textwidth]{v4_crossenc.pdf}
\caption{Distribuição das classificações do \textit{cross-encoder} por estrato. A
fatia laranja na inversão de polo --- pares lidos como \emph{implicação} --- é o
sintoma de que a arquitetura não modela papéis temáticos.}
\label{fig:cross}
\end{figure}

O modo de falha importa mais que a taxa. Dos sete pares de inversão de polo não
detectados como contradição, seis foram classificados como \emph{implicação}:
o modelo julga que ``o autor deve indenizar o réu'' \emph{implica} ``o réu deve
indenizar o autor''. Isto não é ruído estatístico; é a manifestação precisa de
que a inversão de polo não é uma contradição lógico-formal --- ambas as sentenças
podem ser verdadeiras em situações distintas --- e sim uma inversão referencial de
papéis que nenhuma das arquiteturas testadas, nem o bi-encoder por similaridade
nem o \textit{cross-encoder} por inferência, representa adequadamente. A
consequência para o argumento central é direta: a fragilidade da polaridade de
papel não é um artefato da compressão vetorial que um modelo conjunto corrigiria.
É mais profunda, e atravessa também o paradigma de inferência.



O operador poderia reagir dizendo: ``então preciso de um engenheiro que use o
melhor modelo''. A reação é exatamente a errada, e os dados mostram por quê. O
``melhor'' modelo, pela métrica que o operador sabe ler, é o que mais disfarça a
fragilidade. Não se trata de um defeito que uma engenharia melhor elimina, e sim
de uma propriedade da escolha de representar significado por proximidade
distribucional num único vetor --- escolha que é, para os modelos de recuperação
densa que aqui medimos, simultaneamente a fonte de seu poder e a fonte de suas
falhas de polaridade. Cabe delimitar o alcance dessa afirmação. Ela não vale para
modelos generativos autoregressivos, cujo comportamento é moldado também por
camadas de otimização que não se aplicam a um \textit{embedding} --- ajuste por
reforço com preferência humana, alinhamento constitucional, computação inferencial
em tempo de execução --- e cujas falhas têm, por isso, outras fontes além da
geometria distribucional. O objeto deste artigo é o componente de busca, não o
gerador: um \textit{sentence embedding} não possui política de geração, modelo de
recompensa nem cadeia de raciocínio; ele projeta texto em um vetor, e é nessa
projeção que a polaridade se perde. A polaridade referencial, em particular a
inversão de polo, resiste a tudo o que testamos nesse componente: aos quatro
modelos, ao aprendizado contrastivo do e5, e mesmo a uma arquitetura de atenção
cruzada, como a Seção~\ref{sec:cross} documenta em detalhe. O problema não é de
modelo. É de como a linguagem representa papéis e negações, e de como essa
representação sobrevive, ou não, à compressão vetorial.

Nenhuma dessas camadas é visível para quem opera o sistema como um buscador.
Todas são previsíveis para quem entende que está lidando com um modelo estatístico
de linguagem, e agora com a geometria do espaço em que esse modelo vive. A
diferença entre o operador competente e o exposto não é a quantidade de prompts
que conhece. É a presença ou ausência de um modelo mental correto do que a máquina
faz com a linguagem e do que um número de similaridade significa.

% ======================================================================
\section{A escala do dano já documentada}

Quem considera isto teórico não acompanhou os números. Um estudo da Stanford
RegLab sobre alucinação jurídica em modelos de propósito geral encontrou taxas
entre 58\% e 82\%, com desempenho em doutrina substantiva que mal supera o chute
aleatório \cite{dahl2024}. Avaliando ferramentas dedicadas que ancoram respostas
por recuperação aumentada, um segundo estudo ainda registrou 17\% de alucinação no
Lexis+~AI e 33\% no Westlaw \cite{magesh2025}. Recuperação aumentada reduz o
problema; não o elimina. E se a etapa de \textit{embedding} preserva a polaridade
por margem de milésimos, a recuperação que alimenta o gerador pode entregar o
precedente de sentido invertido, sobre o qual o modelo gerativo então constrói uma
resposta fluente e confiante.

E há o caso fundador. Em \emph{Mata v.\ Avianca}, dois advogados protocolaram uma
peça com seis citações geradas por IA; nenhuma existia. Quando um deles pediu ao
modelo que confirmasse se eram reais, ele confirmou \cite{mata2023}. O sistema não
mentiu: gerou a continuação estatisticamente mais provável para ``essas citações
são reais?'', que era a confirmação. O operador leu espelhamento como verificação.
O erro tem a mesma forma do erro de ler cosseno $0{,}99$ como compreensão: em ambos
os casos, uma propriedade estatística do sistema é interpretada como um ato
cognitivo que ele não realizou. Registre-se que a geração de texto envolve camadas
de otimização adicionais --- alinhamento por preferência, raciocínio inferencial ---
que este artigo não mede e cujas contribuições para o erro são objeto de outra
literatura; o ponto aqui é apenas que a \emph{forma} do equívoco do operador, tomar
estatística por cognição, reaparece nos dois extremos do \textit{pipeline}, da
recuperação à geração.

A lição não é ``a IA erra''. É que os erros têm uma gramática, e que ela é também
geométrica. Quem a conhece sabe onde o sistema é frágil: na polaridade, com
severidade máxima na inversão de polo; sob anisotropia, onde a fragilidade se
esconde atrás de escores altos. O operador alfabetizado não confia menos de modo
difuso; desconfia de modo cirúrgico, nos pontos exatos onde a arquitetura e a
geometria preveem a falha.

% ======================================================================
\section{Verificação por raciocínio: a pergunta importa mais que o modelo}
\label{sec:reasoning}

A Seção~\ref{sec:cross} mostrou que a inversão de polo resiste a um
\textit{cross-encoder} de inferência. Cabe perguntar se os modelos generativos de
raciocínio --- a geração que sucede a recuperação em um \textit{pipeline} RAG, e
que a Seção~\ref{sec:limita} delimita como fora do objeto de medição deste artigo
--- conseguem, quando usados como passo explícito de \emph{verificação} sobre o
mesmo fenômeno, capturar a inversão que o \textit{embedding} e o
\textit{cross-encoder} perdem. A questão não é se um modelo de raciocínio
substitui a recuperação, e sim se ele repara, a jusante, a falha de polaridade que
a recuperação deixa passar. O experimento que segue continua, portanto, a
investigação sobre a inversão de polo iniciada na Seção~\ref{sec:resultados}, agora
em uma terceira arquitetura.

Submetemos os 16 pares de inversão de polo a três modelos generativos de uma
mesma família, cobrindo um eixo de capacidade e custo --- de um modelo econômico a
um modelo de fronteira --- sob \emph{duas formulações} da mesma tarefa. A primeira
replica a pergunta do \textit{cross-encoder}: ``estas frases se contradizem?''. A
segunda pergunta pela estrutura argumental: ``a posição das partes foi trocada, ou
cada parte conserva o mesmo papel?''. A distinção é teórica antes de ser empírica:
a inversão de polo não é uma contradição de valor de verdade --- ``o autor deve ao
réu'' e ``o réu deve ao autor'' podem ambas ser verdadeiras em um litígio com
pedidos recíprocos --- e sim uma inversão de papéis temáticos, dimensão ortogonal
à contradição. A hipótese é que a formulação, e não a capacidade do modelo, decida
o resultado.

A Figura~\ref{fig:reasoning} reporta os achados, e eles são mais nítidos que a
hipótese previa. Sob a pergunta por papel temático, todos os três modelos
identificam a inversão na quase totalidade dos casos --- Haiku $13/16$, Sonnet
$15/16$, Opus $16/16$ ---, todos superando com folga o \textit{cross-encoder}
($9/16$). Sob a pergunta por contradição, todos colapsam --- Haiku $7/16$, Sonnet
$0/16$, Opus $3/16$. O salto entre as duas formulações, dentro de cada modelo, é de
seis a quinze pares. A formulação domina.

Mais revelador que as contagens é \emph{quais} pares cada formulação resolve. Os
sete pares em que o \textit{cross-encoder} falhou são os de papéis recíprocos
(apelante/apelada, credor/devedor, requerente/requerido e análogos), em que a
troca produz duas frases logicamente compatíveis. Sob a pergunta por papel
temático, os dois modelos maiores recuperam \emph{todos} os sete ($7/7$), e o
modelo econômico, cinco; o \textit{cross-encoder} recuperara zero. A formulação
correta ataca exatamente a fraqueza estrutural do paradigma de inferência.

\begin{figure}[H]
\centering
\includegraphics[width=0.98\textwidth]{v4_reasoning.pdf}
\caption{Verificação da inversão de polo por modelos de raciocínio. À esquerda:
trocar a pergunta de ``contradição'' para ``papel temático'' eleva a detecção de
quase zero a quase perfeita em todos os modelos, acima do \textit{cross-encoder}
(linha tracejada). À direita: os pares recíprocos que o \textit{cross-encoder}
errou ($0/7$) são recuperados pela pergunta por papel temático.}
\label{fig:reasoning}
\end{figure}

Dois resultados merecem destaque, com a ressalva de que, sobre $16$ itens,
diferenças de poucos pares entre modelos não são estatisticamente significativas;
o que é robusto é o padrão, não o placar.

O primeiro é o domínio da formulação. Que o modelo mais econômico, sob a pergunta
correta, alcance $13/16$ e supere o \textit{cross-encoder} dedicado mostra que a
verificação eficaz de polaridade não requer o modelo de maior capacidade; requer a
pergunta que captura a estrutura argumental. O recurso decisivo não é a potência
da ferramenta, mas a formulação da consulta.

O segundo é contraintuitivo, e as justificativas espontâneas dos modelos o
explicitam. Sob a formulação inadequada, a capacidade não ajuda: o modelo de
fronteira não lidera. Sob ``contradição'', os acertos são Haiku $7$, Sonnet $0$,
Opus $3$ --- a detecção não cresce com a capacidade, e o modelo mais caro fica
abaixo do mais barato. A razão aparece nas respostas. Diante do par ``o autor deve
indenizar o réu'' / ``o réu deve indenizar o autor'', o Opus respondeu que não há
contradição, justificando que, em caso de danos recíprocos, ambos os deveres podem
coexistir e ambas as frases podem ser verdadeiras ao mesmo tempo. O raciocínio é
juridicamente correto sobre valor de verdade --- e é exatamente por ser correto na
dimensão da contradição que falha na dimensão que importa, a do papel das partes.
Mais capacidade aplicada à formulação errada não corrige o erro; produz uma recusa
mais fundamentada da pergunta certa. O modelo econômico, menos cauteloso, por vezes
respondeu ``contradição'' a esses mesmos pares --- acertando o rótulo pela razão
errada, o que não é virtude e não se generaliza.

Um achado lateral, registrado por honestidade, reforça a tese sob outro ângulo. Ao
ser interrogado sobre o par que inverte credor e devedor, o modelo de fronteira
identificou a troca de papéis e, não solicitado, apontou que uma das frases é
juridicamente impossível, pois quem cede um crédito é necessariamente seu titular,
o credor. A observação é correta e expõe uma limitação do próprio conjunto: alguns
pares de inversão de polo, ao trocar mecanicamente os termos, geram uma frase
incoerente em vez de meramente invertida, o que confunde a variável de interesse.
Trata-se de ruído de construção a corrigir em uma versão ampliada do conjunto; não
altera o padrão agregado, mas delimita a precisão dos itens individuais. De modo
análogo, parte dos poucos erros sob a pergunta por papel temático decorreu de os
modelos lerem ``papel'' como função sintática (ambas as partes são sujeito que
sustenta tese) em vez de posição processual concreta --- ambiguidade da própria
pergunta, que ilustra, ela mesma, que a verificação exige formulação precisa.

A consequência estende a tese central do artigo à camada de raciocínio, sem
contrariá-la. Na recuperação, a fragilidade se esconde atrás de um cosseno alto; na
verificação por raciocínio, ela se esconde atrás da formulação plausível mas
inadequada. Em ambos os casos, o que separa o acerto do erro não é o poder do
modelo, e sim a competência de quem o opera --- na recuperação, saber que a
similaridade alta é geometricamente enganosa; na verificação, saber que a pergunta
deve interrogar o papel das partes, não a contradição lógica. A era do raciocínio
não dispensa a alfabetização do operador; desloca-a da geometria do
\textit{embedding} para a estrutura da pergunta de verificação.

% ======================================================================
\section{O que muda quando o operador entende a linguagem da máquina}

Três deslocamentos concretos decorrem da perspectiva adotada. O primeiro, na
verificação: quem entende que o modelo gera a continuação mais plausível para de
pedir ao próprio modelo que confirme os fatos que produziu, e leva a verificação
para a fonte primária. O segundo, na formulação da entrada: compreender o espelho
reativo leva o operador a forçar o contraditório, pedindo a antítese e as objeções
mais fortes à própria posição, em vez de entregar a tese pronta e convidar o
sistema a fabricar suporte. O terceiro, e o que estes resultados tornam urgente,
na arquitetura dos sistemas que o Direito constrói: quem projeta ou contrata um
\textit{pipeline} de recuperação precisa saber que a similaridade absoluta de um
\textit{embedding} é geometricamente enganosa, que a comparação entre modelos deve
ser feita sobre métricas normalizadas e de \textit{ranking}, e que a polaridade ---
sobretudo a inversão de polo --- exige mitigação explícita: \textit{reranking} de
atenção cruzada, tratamento dedicado de pares dialéticos, verificação de
polaridade independente do escore de similaridade. Essas são decisões linguísticas
e geométricas antes de serem decisões de engenharia, e o jurista que não as
entende terceiriza ao fornecedor escolhas que determinam se o sistema sustenta ou
sabota o trabalho.

% ======================================================================
\section{Limitações, escopo e trabalho em andamento}
\label{sec:limita}

Convém ser explícito quanto ao estatuto deste estudo: ele é uma prova de conceito
controlada, não um \textit{benchmark} de escala, e três limitações delimitam o que
ele estabelece e o que fica para a continuação.

\textit{Escala dos dados.} Os 40 pares foram construídos por um único autor com
expertise de domínio. Eles ilustram mecanismos com intervalo de confiança
reportado, e a convergência do padrão entre quatro modelos de gerações e
geometrias distintas sugere robustez; mas não estabelecem distribuições
populacionais, e os estratos de quantificador/modal permanecem pequenos ($n=4$).
A extensão natural é ampliar e submeter os pares a anotação independente.

\textit{Escala da recuperação.} A tarefa de recuperação usa um corpus reduzido
(dois candidatos mais doze distratores). Ela isola a discriminação de polaridade e
mostra que a margem que a sustenta é de milésimos, mas não \emph{prova} o colapso
dessa margem sob ruído de produção. A continuação direta, para a qual o
\textit{script} de replicação já está estruturado, é injetar os pares contrastivos
em um acervo de milhares de \textit{chunks} minerados de jurisprudência pública e
medir a degradação da margem em função do tamanho do corpus. A hipótese ---
sustentada pela estreiteza já observada --- é que distratores reais estreitam, não
alargam, a margem; convertê-la em evidência quantitativa é trabalho em andamento.

\textit{Modelos proprietários.} Medimos quatro modelos abertos, um deles de
fronteira. A espinha dorsal de parte do mercado jurídico são \textit{embeddings}
proprietários (OpenAI, Cohere), que empregam técnicas próprias de redução de
dimensionalidade. Sua geometria pode diferir, e testá-los é necessário para
generalizar a afirmação ao cenário comercial completo. Registramos isto como
limitação franca: a tese de que a anisotropia mascara a fragilidade está
estabelecida para o modelo aberto de fronteira e fundamentada na literatura sobre
o fenômeno, mas sua confirmação nos modelos fechados de maior uso comercial fica
como passo seguinte.

\textit{Método de isotropização.} Normalizar pela anisotropia é uma entre várias
técnicas. Escolhemos a normalização por \textit{baseline} de pares mínimos por ser
a mais transparente e diretamente interpretável \cite{heros2023}; uma validação
cruzada por remoção de componentes principais (\textit{whitening}) é caminho
complementar para confirmar que o reaparecimento da estrutura não depende do
método particular.

\textit{Validade dos pares de inversão de polo e robustez ao subconjunto limpo.}
A inversão de polo sustenta a afirmação mais forte do artigo --- a de que é o
estrato mais severo. Isso exige escrutínio dos próprios itens, pois inverter
mecanicamente os papéis pode gerar não uma frase invertida, mas uma frase
juridicamente impossível. Auditamos manualmente os $16$ pares: $12$ são
inequivocamente coerentes (a inversão produz duas frases ambas possíveis em
situações distintas), $3$ são limítrofes (ambíguos quanto a legitimidade ativa
\textit{versus} passiva, mas não impossíveis) e $1$ é impossível --- o par que
inverte credor e devedor, pois apenas o titular do crédito pode cedê-lo, defeito
identificado espontaneamente pelo próprio modelo de fronteira. Reexecutamos as
análises de que dispomos sobre o subconjunto de $12$ pares limpos, e o padrão não
apenas sobrevive como se reforça. O \textit{cross-encoder} continua a falhar na
mesma proporção (detecta contradição em $7/12$, contra $9/16$ no conjunto cheio), e
a verificação por papel temático permanece quase perfeita (Opus e Sonnet $12/12$,
Haiku $10/12$). Quanto à métrica geométrica, remedimos o cosseno normalizado da
inversão de polo no e5 sobre os três conjuntos, e o resultado é o achado mais forte
do artigo: o \textit{gap} normalizado entre a paráfrase de controle e a inversão de
polo é \emph{negativo} --- $-0{,}016$ nos $16$ pares e $-0{,}034$ nos $12$ limpos.
Um \textit{gap} negativo significa que o e5 atribui à inversão de polo um cosseno
normalizado ($0{,}97$ no subconjunto limpo) \emph{superior} ao que atribui a duas
paráfrases legítimas ($0{,}94$): o modelo considera ``o autor deve ao réu'' e ``o
réu deve ao autor'' mais próximas entre si do que dois sinônimos verdadeiros. E o
efeito é \emph{mais} forte nos pares limpos --- os pares imperfeitos atenuavam-no,
de modo que a depuração reforça, em vez de comprometer, a afirmação de que a
inversão de polo é o estrato mais severo. A contaminação existia, mas operava
contra a tese, não a favor; removê-la fortalece o achado. (Distinga-se este
\textit{gap} da inversão de polo isolada, negativo, do \textit{gap} médio de todos
os estratos contrastivos reportado na Seção~\ref{sec:aniso}, de $0{,}09$; são
granularidades distintas, e a inversão de polo é o estrato que puxa a fragilidade
ao extremo.) O conjunto será ainda depurado e ampliado, com os pares limítrofes
reescritos para isolar legitimidade ativa de passiva, em versão futura.

\textit{Verificação por raciocínio.} O experimento da Seção~\ref{sec:reasoning}
usa esses mesmos pares e três modelos de uma única família; suas contagens
individuais estão sujeitas ao $n$ pequeno. O padrão central --- domínio da
formulação sobre a capacidade e recuperação dos recíprocos --- é robusto, inclusive
ao subconjunto limpo acima; as diferenças de poucos pares entre modelos não o são,
e não são reivindicadas.

Nenhuma dessas limitações enfraquece a tese central, que diz respeito à
\emph{interpretação} dos escores de similaridade e foi confirmada em todos os
modelos medidos. Elas delimitam a \emph{generalização} da evidência, e o caminho
para ampliá-la está especificado e instrumentado.

% ======================================================================
\section{Conclusão}

A tese é esta. O operador do Direito não precisa abordar a operação de IA pela
perspectiva linguística porque ela é elegante, mas porque o objeto que opera é, na
sua natureza mais profunda, um artefato linguístico e geométrico, e as falhas que
ele produz têm a forma da representação distribucional do significado em um espaço
vetorial enviesado. Este artigo tornou uma dessas falhas mensurável em português e
mostrou algo que a intuição corrente inverte: o modelo de busca mais avançado não
é mais seguro contra o erro de polaridade --- é mais perigoso, porque o esconde
melhor. A inversão de polo resiste a quatro gerações de modelos e a uma arquitetura
de inferência; o aprendizado contrastivo do modelo de fronteira não a removeu,
apenas a deslocou para um regime anisotrópico onde um cosseno de $0{,}99$ disfarça
uma discriminação de cinco milésimos. O operador que lê esse número como certeza
comete o erro exato que a competência jurídica não o ensina a ver.

A competência jurídica diz ao operador se um precedente é aplicável. Não lhe diz
se a máquina o está recuperando porque o sustenta ou porque o contradiz, nem com
que margem, nem se o escore de similaridade que ele exibe é discriminação real ou
artefato da geometria do espaço. Para responder a essas perguntas, que decidem se
a IA será instrumento ou armadilha, é preciso uma alfabetização que a formação
jurídica tradicional não fornece e que nenhum manual de atalhos substitui. É
preciso entender a linguagem da máquina --- e a geometria em que ela vive.

% ======================================================================
\paragraph{Disponibilidade de material.} O conjunto de pares mínimos e os scripts
de medição (cosseno, normalização por anisotropia, \textit{cross-encoder} e
recuperação) são reproduzíveis com \texttt{sentence-transformers} em CPU. Os
resultados do modelo de fronteira foram obtidos com \texttt{multilingual-e5-large}.

\paragraph{Declaração sobre uso de IA.} A redação foi assistida por ferramentas de
IA sob revisão integral do autor. Os experimentos, dados e interpretação são de
responsabilidade do autor.

% ======================================================================
\begin{thebibliography}{9}
\small

\bibitem{dahl2024}
DAHL, M.; MAGESH, V.; SUZGUN, M.; HO, D.~E.
\emph{Large Legal Fictions: Profiling Legal Hallucinations in Large Language
Models}. Journal of Legal Analysis, v.~16, n.~1, p.~64--93, 2024.
\url{https://doi.org/10.1093/jla/laae003}.

\bibitem{ethayarajh2019}
ETHAYARAJH, K.
\emph{How Contextual are Contextualized Word Representations? Comparing the
Geometry of BERT, ELMo, and GPT-2 Embeddings}. In: Proceedings of EMNLP-IJCNLP,
2019, p.~55--65. \url{https://doi.org/10.18653/v1/D19-1006}.

\bibitem{hammerl2023}
HÄMMERL, K.; LIBOVICKÝ, J.; FRASER, A.
\emph{Exploring Anisotropy and Outliers in Multilingual Language Models for
Cross-Lingual Semantic Sentence Similarity}. In: Findings of ACL, 2023.
\url{https://doi.org/10.18653/v1/2023.findings-acl.439}.

\bibitem{harris1954}
HARRIS, Z.~S.
\emph{Distributional Structure}. Word, v.~10, n.~2--3, p.~146--162, 1954.
\url{https://doi.org/10.1080/00437956.1954.11659520}.

\bibitem{heros2023}
LIN, Y.-T.; CHEN, Y.-N.
\emph{HEROS: A Diagnostic Probe of Sentence Encoders via Minimal Pairs} (HEROS).
Preprint, 2023. \url{https://arxiv.org/abs/2306.05083}.

\bibitem{magesh2025}
MAGESH, V.; SURANI, F.; DAHL, M.; SUZGUN, M.; MANNING, C.~D.; HO, D.~E.
\emph{Hallucination-Free? Assessing the Reliability of Leading AI Legal Research
Tools}. Journal of Empirical Legal Studies, 2025.
\url{https://doi.org/10.1111/jels.12413}.

\bibitem{mata2023}
ESTADOS UNIDOS. United States District Court for the Southern District of New
York. \emph{Mata v.\ Avianca, Inc.}, No. 22-cv-1461 (PKC). Opinion and Order on
Sanctions, 22 jun. 2023.

\bibitem{reimers2019}
REIMERS, N.; GUREVYCH, I.
\emph{Sentence-BERT: Sentence Embeddings using Siamese BERT-Networks}. In:
Proceedings of EMNLP-IJCNLP, 2019, p.~3982--3992.
\url{https://doi.org/10.18653/v1/D19-1410}.

\end{thebibliography}

\end{document}
